% ======================================================
%  Zagadki -- poziom zaawansowany. Hand-written, for the reason the basic
%  volume's are: a word trap is a joke and a joke has an author.
%
%  These are harder than the basic volume's in a specific way rather than a
%  general one -- each has a short answer that is easy to check and a route to
%  it that is not the first one anybody tries.
% ======================================================

\begin{zestaw}{Zagadki — trudniejsze}{3}{20:00}{1}
  \zz{Trzy skrzynki: „jabłka”, „pomarańcze”, „mieszane”. \;/\; Każda etykieta jest błędna. \;/\; Ile owoców trzeba wyjąć, żeby ustalić zawartość wszystkich trzech?}{1 \hfill {\itshape\footnotesize Ze skrzynki „mieszane” — skoro etykieta kłamie, jest jednorodna.}}
  \zz{Dwa sznurki, każdy pali się dokładnie godzinę, \;/\; ale nierównomiernie. \;/\; Jak odmierzyć 45 minut?}{Zapal I z obu końców i II z jednego \hfill {\itshape\footnotesize I gaśnie po 30 min; wtedy podpal drugi koniec II — zostanie mu 15.}}
  \zz{Osiem monet, jedna cięższa, waga szalkowa bez odważników. \;/\; Ile ważeń wystarczy w najgorszym razie?}{2 \hfill {\itshape\footnotesize 3 na 3, potem 1 na 1 — trzy grupy, nie dwie.}}
  \zz{Dwanaście monet, jedna inna — nie wiadomo, cięższa czy lżejsza. \;/\; Ile ważeń wystarczy?}{3 \hfill {\itshape\footnotesize Trzy ważenia rozróżniają $3^3 = 27$ przypadków; tu jest ich 24.}}
  \zz{Cztery osoby przechodzą nocą przez most, jedna latarka, \;/\; idą po dwie, w tempie wolniejszego. \;/\; Czasy: 1, 2, 5, 10 minut. Najkrótszy łączny czas?}{17 minut \hfill {\itshape\footnotesize Dwaj najwolniejsi idą razem — raz.}}
  \zz{Dzbanki na 3 i 5 litrów, kran bez podziałki. \;/\; Jak odmierzyć dokładnie 4 litry?}{Napełnij 5, przelej do 3, zostaną 2 \hfill {\itshape\footnotesize Opróżnij trójkę, wlej te 2, dopełnij piątkę i dolej 1 do trójki.}}
  \zz{Ciąg: $1$, $11$, $21$, $1211$, $111221$, \ldots \;/\; Jaki jest następny wyraz?}{312211 \hfill {\itshape\footnotesize Każdy wyraz opisuje poprzedni: „trzy jedynki, dwie dwójki, jedna jedynka”.}}
  \zz{Iloma zerami kończy się $100!$?}{24 \hfill {\itshape\footnotesize Tyle, ile piątek w rozkładzie: $20 + 4$.}}
  \zz{Ile liczb od 1 do 1000 zawiera cyfrę 7?}{271 \hfill {\itshape\footnotesize Łatwiej policzyć te bez siódemki: $9 \times 9 \times 9 = 729$.}}
  \zz{Ile razy na dobę wskazówki zegara \;/\; tworzą kąt prosty?}{44 \hfill {\itshape\footnotesize 22 razy na 12 godzin — nie 24, bo nakładanie się „zjada” dwa.}}
  \zz{Mam tyle samo braci co sióstr. \;/\; Moja siostra ma dwa razy więcej braci niż sióstr. \;/\; Ilu nas jest?}{7 \hfill {\itshape\footnotesize 4 braci i 3 siostry; ja jestem bratem.}}
  \zz{Sznur opina Ziemię po równiku. \;/\; Dokładasz do niego 1 metr. \;/\; Jak wysoko uniesie się nad powierzchnię?}{Około 16 cm \hfill {\itshape\footnotesize $1/(2\pi)$ — promień Ziemi się skraca.}}
  \zz{Siatka $3 \times 3$ kwadratów. \;/\; Ile jest dróg z lewego górnego rogu do prawego dolnego, \;/\; jeśli wolno iść tylko w prawo i w dół?}{20 \hfill {\itshape\footnotesize Sześć kroków, z czego trzy w prawo: $C(6,3)$.}}
  \zz{Pierwszą połowę trasy jedziesz 60 km/h. \;/\; Jak szybko musisz jechać drugą połowę, \;/\; żeby średnia wyszła 120 km/h?}{Nie da się \hfill {\itshape\footnotesize Na samą pierwszą połowę zużyłeś już cały czas przewidziany na całość.}}
  \zz{W pokoju jest $n$ osób. \;/\; Przy jakim najmniejszym $n$ szansa, że dwie mają urodziny tego samego dnia, \;/\; przekracza połowę?}{23 \hfill {\itshape\footnotesize Liczy się liczba par, a nie osób: $C(23,2) = 253$.}}
  \zz{Kartka złożona na pół 50 razy \;/\; miałaby grubość rzędu?}{Odległości do Słońca \hfill {\itshape\footnotesize $2^{50}$ razy $0{,}1$ mm to około $10^{8}$ km.}}
\end{zestaw}
